\chapter{选题报告}\label{chap:selection-report}

\section{选题依据}

\subsection{研究来源}

本课题来源于前期Agent运行可观测原型实践与相关文献研究。在原型实践中，运行日志、时间线和结构视图已经能够呈现模型消息、工具调用、请求信息与部分上下文来源，但开发者仍需在多个视图之间自行判断某一运行现象由哪个Harness机制塑造、当前判断有哪些证据支持。该问题构成了本研究的直接来源。

\subsection{研究背景}

大语言模型Agent正由单轮问答工具转向能够调用外部工具、访问工作环境、维护上下文并持续完成复杂任务的复合系统。此类系统的最终行为并不只由基础模型决定，还受到执行环境、工具接口、上下文与记忆、生命周期编排、可观测、验证评估以及治理安全等运行基础设施的共同影响。围绕模型并负责把一次任务组织成可执行过程的这组基础设施，通常被称为Agent Harness。随着Agent任务链条延长，开发者面对的核心困难也从“结果是否正确”进一步转向“为什么形成这一结果，以及下一步应检查哪个机制”。

现有评价工具能够给出任务成功率、基准得分、成本与延迟，却难以把结果直接转化为可操作的系统修改。针对生产环境中LLM产品实践者的访谈显示，19名受访者中有17名经历过“结果—可行动性鸿沟”：团队能够收集分数和反馈，却难以判断问题来自提示组织、检索、工具、上下文还是其他系统部件\cite{ref01}。这一现象在Agent场景中更加突出。一次失败轨迹往往同时包含自然语言推理、工具调用、环境反馈和中间状态变化，异常可能在早期产生、在后续步骤传播，并最终以错误结果或提前终止的形式出现。

技术研究已经开始把原始轨迹组织为数据流、控制流和实现锚点，并据此进行自动归因与Harness修复。HarnessFix通过Harness-aware Trace Intermediate Representation把运行步骤与提示模板、工具规范、配置、验证脚本等实现工件对齐，在四类基准上取得6.3\%至18.4\%的性能提升\cite{ref20}。这些成果说明轨迹—机制关联具有明确的工程价值，但其主要消费者仍是自动诊断与修复模块。自动系统能够生成关联或候选原因，并不等同于开发者能够正确理解这些关系。已有自动故障归因研究中，责任Agent的识别准确率约为53.5\%，具体步骤的准确率仅为14.2\%\cite{ref03}，表明系统输出本身仍需要核查。

另一方面，Agent可观测与交互调试工具已能提供时间线、层级视图、断点、回放和现场编辑。AGDebugger、AgentStepper、AgentGUI等研究证明，结构化轨迹有助于浏览长会话、识别部分缺陷并降低操作负担\cite{ref02,ref13,ref18}。然而，现有界面常把原始记录、规则推导和模型生成的候选解释放在相近的视觉层级，用户容易把“界面画出的关系”理解为“系统已经证明的因果关系”。例如，工具结果出现在日志中，只能证明该结果产生过；若缺少后续请求快照，就不能断言它已被模型消费。固定结构图中的连线也可能表示预设架构，而非本次运行真实经过的路径。

因此，本研究把问题限定为Agent Harness的运行理解：开发者如何把一次运行中的现象与背后起作用的Harness机制联系起来，并知道自己凭什么可以这样联系。研究进一步聚焦上下文构建与变换过程，考察信息从何处进入模型请求、在后续步骤中如何被保留、替换、裁剪或摘要，以及现有证据能否支持某项故障判断。相较于覆盖全部Harness层，这一范围既保留了问题的代表性，也便于构造具有客观真值的研究材料。

\subsection{研究目的及意义}

本研究旨在构建一套面向开发者运行理解的轨迹—机制关联与证据状态呈现方法。具体而言，首先通过真实运行案例和形成性研究识别开发者在上下文来源追踪、信息变换追踪和故障定位中的证据需求与误判模式；其次在既有Architecture原型基础上，将数据流、控制流与实现锚点组织为可下钻的关联视图，并把每项关系区分为“实际记录到”“按规则推导”和“目前无法知道”三类证据状态；最后通过具有客观真值的三条件受控实验，检验关联组织与证据状态表达对理解正确率、证据定位、完成时间和无依据结论的影响。

本研究的理论意义主要体现在三个方面。第一，将传统程序调试中的why与why-not问题、信息觅食和层级排故研究延伸到由自然语言推理、工具交互与运行时上下文共同构成的Agent系统，检验既有调试认知理论在概率性运行环境中的适用边界。第二，引入Data–Frame意义建构框架，解释开发者如何在轨迹证据与机制假设之间循环修正理解，并把“无依据结论”作为运行理解质量的重要表现。第三，把证据状态作为Agent可观测界面的独立设计变量，讨论记录、推导与未知如何影响使用者对系统结论的校准。

本研究的实践意义在于为Agent可观测和调试工具提供可检验的交互设计依据。研究不以继续增加日志字段为目标，而是关注已有证据如何被组织、核验和限定。预期成果包括一组可复用的运行理解任务、一个证据状态显式的轨迹—机制关联原型、一套受控评价方法以及相应设计原则。上述成果可为Agent开发平台、可观测工具和自动诊断系统的人机接口提供参考，使开发者能够更快定位证据，同时在证据不足时保留“尚不能判断”。

\section{国内外研究综述}
本报告的国内外研究综述见\hyperref[chap:literature-review]{第\ref{chap:literature-review}章“文献综述”}，此处不再重复排版。

\section{研究问题与方法}

\subsection{研究问题}

本研究关注一次Agent运行已经发生、系统暂不自动修复的情境。核心问题不是简单增加日志数量，而是帮助开发者把分散的运行记录组织为可以检查、可以说明依据且能够保留未知边界的机制理解。基于前述研究现状与不足，提出以下三个研究问题：

RQ1：开发者在理解Agent Harness运行过程和定位故障时，需要哪些跨步骤、跨组件的证据关联？现有的原始轨迹与静态结构视图会在哪些任务上导致找不到证据、误解连线或把推测当作事实？

RQ2：在底层运行记录相同的条件下，把分散事件组织为数据流、控制流与实现锚点的轨迹—机制关联，能否提高开发者的运行理解正确率、证据定位准确率并缩短定位时间？

RQ3：在轨迹—机制关联已经呈现的基础上，进一步标明“实际记录到”“按规则推导”和“目前无法知道”三类证据状态，能否减少开发者在证据不足时作出确定判断的比例，并在认知负荷可接受的前提下改善结论校准？

\subsection{研究方法}

本研究采用“形成性研究—原型设计—受控评估”的混合研究路径。研究一识别实际困难并形成理解任务集，研究二把困难转化为设计要求并实现原型，研究三通过条件隔离检验界面组织与证据状态表达的作用。三个阶段使用同一组核心概念和可追溯的材料链，保证研究发现能够对应到设计机制，设计机制又能够被实验评价。

\subsubsection{第一阶段：形成性研究}

第一阶段拟招募8至12名具有Agent开发、调试或高频使用经验的参与者。研究材料由预先采集的真实运行记录和少量经过仪器化处理的构造案例组成，重点覆盖三类情境：上下文来源混淆，即某条约束、历史信息或检索材料从何处进入当前请求；压缩、裁剪或替换后的信息变化，即某条内容在后续请求中是否仍然存在；工具结果进入后续请求的过程，即工具输出是否被保存、转换并被模型实际看见。

参与者使用现有的原始轨迹和静态Architecture视图完成排查任务，并在操作过程中出声思考。研究记录其导航路径、视图切换、证据选择、假设变化、放弃点和最终结论，任务后进行半结构化访谈。分析采用反思性主题分析，并以Data–Frame理论中的框架形成、扩展、质疑、比较和重构活动作为敏化概念，同时允许新的困难类型从材料中涌现。阶段性产出包括困难清单、理解任务集、证据需求清单和初步设计要求。

为避免把“界面更直观”误当作研究贡献，本阶段设置明确的继续条件：至少识别出能够在多个参与者或多个案例中重复出现、并实际影响答案正确性的困难模式；若主要反馈仅涉及视觉偏好而没有稳定的证据困难，则缩小研究范围或重新构造任务，不直接进入正式原型评估。

\subsubsection{第二阶段：关联界面原型设计}

第二阶段以已有Architecture原型为基础进行迭代。原型保留固定运行结构、会话—轮次—步骤联动以及上下文构成与运行链路并列三项特征，并新增三类核心机制。

第一，建立轨迹事件与Harness机制之间的关联。数据流关系用于表示信息的来源、传递、变换与消失；控制流关系用于表示重试、委托、校验、终止等运行决策；实现锚点用于把运行时证据连接到提示模板、工具规范、配置项、上下文处理函数或验证脚本。第二，为每项关联增加证据状态。实际记录到的关系必须能够返回原始事件；按规则推导的关系必须能够展开其输入、规则和适用前提；目前无法知道的关系需说明缺少什么信号以及可能的补证方式。第三，支持从症状到机制和从系统结构到事件的双向下钻，以兼容症状搜索与拓扑搜索两种排故策略。

原型开发遵循“发现—设计要求—界面机制”的映射。每一项核心交互都需标明其来自研究一的哪类困难，并在研究三中对应至少一个客观指标。原型将先经过专家走查和小规模可用性测试，修正术语理解、状态识别、跨视图选中和返回原始证据等问题，再进入正式实验。

\subsubsection{第三阶段：三条件受控实验}

研究三采用三条件对照设计。所有条件使用相同的运行案例和底层记录，仅改变信息组织与证据状态表达，以避免把“增加了更多数据”和“设计得更好”混为一谈。

\begin{table}[htbp]
\centering
\caption{实验条件设置}
\label{tab:report-7}
\scriptsize
\renewcommand{\arraystretch}{1.18}
\begin{tabularx}{\textwidth}{>{\raggedright\arraybackslash}p{2.7cm} *{2}{>{\raggedright\arraybackslash}X}}
\toprule
\textbf{条件} & \textbf{呈现方式} & \textbf{主要比较目的} \\
\midrule
A 常规轨迹 & 以时间顺序呈现模型请求、工具调用、结果和状态变化 & 作为现有日志查看方式的基线 \\
B 轨迹—机制关联 & 在A的记录基础上增加数据流、控制流和实现锚点，并支持双向下钻 & 与A比较，检验关联组织的作用 \\
C 关联加证据状态 & 在B的基础上明确标注记录、推导与未知，并提供状态依据 & 与B比较，检验证据状态表达的作用 \\
\bottomrule
\end{tabularx}
\end{table}

实验任务覆盖事件顺序理解、信息来源追踪、上下文变换判断、机制定位和证据充分性判断。为降低学习效应，条件顺序采用拉丁方平衡，正式任务前设置不计分的训练案例；不同条件使用结构等价但内容不同的材料。参与者背景将记录Agent开发经验、传统调试经验与可视化使用经验，用于检验个体差异。

\subsubsection{材料构建与客观真值}

运行理解研究必须具备可核对的标准答案，否则只能评价“是否觉得清楚”。本研究采用“真实运行记录加仪器化构造案例”的组合。真实案例用于保留长轨迹、噪声和非理想数据等生态特征；构造案例用于有计划地操纵上下文保留、替换、裁剪、摘要和工具结果注入等机制，使关键问题具有明确答案。

每个案例均建立独立的证据清单，记录原始事件、稳定标识、关联规则、实现位置和不可观测边界。案例由两名研究者依据统一手册独立标注，并报告评分者间一致性；存在分歧的案例经讨论后修订或剔除。对“未知”状态的判定同样需要依据：只有在必需信号确实未记录、采样或超出观测边界时，才标记为目前无法知道。该过程既保证实验真值，也为原型中的状态说明提供数据基础。

\subsubsection{测量与数据分析}

研究三的主要因变量包括运行理解正确率、证据定位准确率、完成时间和无证据下结论率。运行理解正确率衡量参与者对事件来源、变换机制和故障位置的判断与客观真值是否一致；证据定位准确率衡量其能否找到支撑结论的轨迹片段与实现锚点；无证据下结论率用于统计参与者在证据不足时仍作肯定判断，或把规则推导误称为直接事实的比例。辅助指标包括导航路径、视图切换、答案修改、Raw NASA-TLX与主观可用性评价。

\begin{table}[htbp]
\centering
\caption{实验测量维度与指标}
\label{tab:report-8}
\scriptsize
\renewcommand{\arraystretch}{1.18}
\begin{tabularx}{\textwidth}{>{\raggedright\arraybackslash}p{2.7cm} *{2}{>{\raggedright\arraybackslash}X}}
\toprule
\textbf{测量维度} & \textbf{主要指标} & \textbf{评价重点} \\
\midrule
理解结果 & 运行理解正确率、故障定位正确率 & 是否形成与真实机制一致的解释 \\
证据能力 & 证据定位准确率、关键证据召回率 & 是否找到并引用了正确证据 \\
结论校准 & 无证据下结论率、未知保留率 & 是否区分事实、推导与未知 \\
效率与负荷 & 完成时间、导航步数、Raw NASA-TLX & 设计收益是否伴随过高成本 \\
\bottomrule
\end{tabularx}
\end{table}

量化数据在试次层面采用混合效应模型分析，以界面条件和任务类型为固定效应，以参与者和材料为随机效应；二值结果使用广义线性混合模型，连续结果根据分布选择线性混合模型或稳健方法。重点比较A与B、B与C，而不是仅报告三组总体差异。样本量由预测试所得效应和随机效应方差通过模拟功效分析确定，不预先沿用旧方案的固定人数。交互日志、屏幕记录和访谈用于解释“界面没有提供线索”“参与者没有发现线索”以及“发现后仍作出错误推断”三类机制差异。

\subsubsection{可行性与伦理}

本研究已具备运行记录、文献基础和Architecture原型，可在不搭建在线自动修复平台的条件下开展。材料使用公开框架、自有原型和预先采集的运行记录，不依赖企业内部数据。涉及人的研究将在取得伦理审查批准后实施，参与者均需完成知情同意；数据以研究编号替代身份信息，屏幕记录和访谈转录在去标识后分析，参与者可以随时退出。研究三将在正式收集数据前预注册研究设计、主要指标与分析计划，并在不泄露个人信息的前提下开放任务材料、评分手册和分析代码。

\section{研究框架}
\begin{figure}[htbp]
\centering
\includegraphics[width=0.92\textwidth]{source-image2.png}
\caption{研究框架与实施路径}
\label{fig:research-framework}
\end{figure}


\section{计划进度}

本研究计划以形成性研究、原型设计和受控评估为主线推进。案例真值构建、伦理审查和参与者招募将提前并行开展，以降低后续实验等待成本。各阶段核心任务和预期成果如下。

\begin{table}[htbp]
\centering
\caption{研究计划与进度安排}
\label{tab:report-9}
\scriptsize
\renewcommand{\arraystretch}{1.18}
\begin{tabularx}{\textwidth}{>{\raggedright\arraybackslash}p{2.6cm} *{3}{>{\raggedright\arraybackslash}X}}
\toprule
\textbf{阶段} & \textbf{时间} & \textbf{核心任务} & \textbf{阶段成果} \\
\midrule
第一阶段：准备与案例构建 & 2026.09—2026.10 & 完成选题报告与伦理申请；核验原型数据边界；采集真实运行记录；建立案例真值与标注手册 & 材料库v0；伦理申请；开题答辩 \\
第二阶段：形成性研究 & 2026.11—2027.01 & 招募8—12名参与者；完成案例排查、出声思考与访谈；归纳困难模式和证据需求 & 困难清单；理解任务集；设计要求 \\
第三阶段：原型设计与预测试 & 2027.02—2027.06 & 迭代轨迹—机制关联与证据状态界面；开展专家走查和可用性预测试；修订实验材料 & 可交互原型；材料库v1；预测试报告 \\
第四阶段：受控实验 & 2027.07—2027.10 & 完成功效分析与预注册；实施三条件实验；整理交互日志、量表与访谈数据 & 正式实验数据集；预注册文档 \\
第五阶段：分析与论文撰写 & 2027.11—2028.02 & 开展混合效应分析和过程分析；形成设计原则；撰写学位论文与阶段论文 & 研究结果；设计原则；论文初稿 \\
第六阶段：修改与答辩 & 2028.03—2028.05 & 完成论文修改、查重、送审和答辩准备；整理开放材料与原型说明 & 学位论文定稿；答辩材料；开放材料 \\
\bottomrule
\end{tabularx}
\end{table}

